%======================================================================
\section{Discussion}\label{sec:disc}
%======================================================================

\paragraph{What the correction is for.}
On the three benchmarks the walk correction does not change top-$K$
accuracy, and in simulation the uncorrected DR operator can even agree
slightly better with the full-exposure ranking, because its bias adds weight
to the imputation (Section~\ref{sec:mc}). The correction matters wherever the
\emph{value} of a propagated score is used rather than only its order within
one user's list: when scores are compared across users, as in the
exposure-capped allocation; when they are aggregated into an estimate of
full-exposure utility; and when an auditor wants to verify that expected
scores do not depend on exposure. In these uses the uncorrected estimate is
biased by 1.8 to 15 times the target over all pairs in the model of
Table~\ref{tab:mc} (0.09 to 2.2 times in the models of Table~\ref{tab:protocol}), and the
bias grows with the number of hops and with small propensities. Two limits
apply. On the unexposed candidates that are actually recommended, DRUP
estimates the full-exposure score only up to the imputation error at the
candidate (Corollary~\ref{cor:cand}), which is 12\% of the mean target at three
hops in our simulation. And the guarantees assume nuisances that do not depend
on the log: with nuisances cross-fitted on the same log, the corrected
operator is not less biased than the uncorrected one in simulation
(Section~\ref{sec:mc}), so a user of the score values needs an independent
log or the sample split of DRUP-split. We have no evidence that the correction changes
rankings, within users or under exposure caps. A practitioner who only needs
rankings loses nothing by using the uncorrected DR adjacency; one who uses
score values across users or pairs, audits exposure invariance, or propagates
over five hops or more, should use the corrected operator. In both cases, most of
the accuracy comes from training-free propagation over a doubly robust graph.

\paragraph{Stakeholders.}
For \emph{item providers}, exposure invariance means that an item is not
penalised in expectation for having been shown rarely in the past. It does not
mean equal exposure: when quality is concentrated, DRUP's lists can be more
concentrated than those of the logged graph (Table~\ref{tab:fair}), and the
scores of rarely exposed items have a larger variance
(Theorem~\ref{thm:var}(a)), so their realised rankings are noisier. Providers
who need exposure guarantees need an explicit mechanism such as the
exposure-capped allocation. For \emph{users}, the attributions show which of
their own interactions drive a recommendation, or certify that none can
overturn it; the population-level part of the score
remains unexplained (it decides 94\% of the top recommendations on Coat and 52\% on KuaiRec), and explaining it through $G$ would conflict with the
privacy guarantee. For the \emph{platform}, DRUP is cheap and training-free
but needs propensities, and the clip $\tau$ and the control-variate weight
$\lambda$ trade exposure invariance against accuracy. This trade-off is a policy
choice rather than a tuning detail. One option is to select the most accurate
configuration subject to a bound on the elasticity measured in a thinning
audit such as ours. For \emph{auditors and regulators}, DRUP produces concrete
artefacts: a single item operator, a thinning audit with a known answer, and
certificates. Under the EU Digital Services Act~\cite{dsa2022}, platforms must
explain the main parameters of their recommender systems (Art.~27), assess and
mitigate systemic risks (Arts.~34--35), undergo independent audits (Art.~37)
and give vetted researchers access to data (Art.~40). Our artefacts fit the
last three, but the audit of exposure invariance needs the platform's
propensities, which an external auditor usually does not have.

\paragraph{Limitations.}
The analysis is static. It assumes a single logging policy and fixed potential
outcomes, whereas recommenders are retrained on the feedback they
generate~\cite{chaney2018confounding,hazrati2022evolution} and can change the
preferences they predict~\cite{perdomo2020performative}; extending the
guarantees to adaptive and slate logging~\cite{swaminathan2017slate} is open.
Proposition~\ref{prop:dep} covers dependence within a user's history, not
dependence across users. The idempotent correction relies on binary outcomes;
graded relevance would need separate unbiased estimates of the powers of $Y$.
Items that were never exposed fall outside $\pi p\ge\tau$ and are scored by the
imputation only. In our main protocol the nuisances are cross-fitted on the
same log, the Yahoo!\,R3 configuration takes degrees from edge estimates, and
the score constants are computed from the log, so the theorems do not cover
these experiments (Remark~\ref{rem:xfit}); the sample split restores them at
the price of an imputation bias on the split-off pairs and, on sparse logs,
of accuracy. The significance tests are conditional
on the single log and its fitted nuisances, and the methods searched grids of
different sizes, which we analyse but cannot fully equalise
(Section~\ref{sec:acc}). The joint-DP guarantee requires public nuisances, and on the
small Coat data the public-nuisance variant loses much of its accuracy. The
certificates are informative only against a handful of fake profiles (on
KuaiRec, 95\% of users at five profiles, none at fifty), and the variance
and concentration bounds are loose by several orders of magnitude
(Section~\ref{sec:mc}). The target $F^*$ is a linear LightGCN score of the
full-exposure world, a normative choice that other stakeholders may not
share. Like every method in our comparison, DRUP selects its hyper-parameters
on a small unbiased sample; without such a sample, selection would need an
off-policy estimate of the validation metric. Finally, our largest catalogue
has about $10^4$ items.

%======================================================================
\section{Implications}
%======================================================================

\paragraph{For research.}
Propagation-level debiasing needs to consider the propagation operator itself,
not only the edges it propagates: multiplying unbiased edge estimates does not
give an unbiased multi-hop score. The repeated-walk bias is exactly
removable, and the tools that remove it (idempotent reduction and Möbius
inversion over walk-index coincidences) apply to any polynomial graph filter
built from inverse-propensity or DR estimates. Our evaluation also shows that
training-free propagation over a DR graph is a strong baseline against which
trained debiased graph recommenders should be compared, and that claims of
exposure invariance can be tested on real data by thinning exposures with a
known factor.

\paragraph{For practice.}
At three hops DRUP runs at the asymptotic cost of one item Gram matrix and needs no gradient training,
so it can serve as an auditable baseline or as a candidate scorer. The
choice between the corrected and the uncorrected operator depends on whether
score values or only within-user orders matter; score values need nuisances
that do not depend on the log, such as those of DRUP-split, since with
nuisances cross-fitted on the same log they carry no unbiasedness guarantee.
On sparse logs, exact exposure
invariance requires small clips and costs accuracy (Table~\ref{tab:tau}). The
structural properties (exact attributions, certificates, a private item
operator and capped allocation) come with any linear propagation operator and
hold for the model with its nuisances held fixed; platforms should report them
together with that scope.

%======================================================================
\section{Conclusion}
%======================================================================

Inverse-propensity and doubly robust graph propagation are biased beyond one
hop because walks that repeat an edge multiply its estimate with itself. We
gave the exact bias, whose worst case grows as the inverse of the propensity
clip, and removed it with an idempotent correction that is exact for any odd
number of hops and, at three hops, has the asymptotic cost of the uncorrected
operator.
With nuisances that do not depend on the log, for instance from a sample
split, the resulting estimator, DRUP, is edge-wise doubly robust for
propagation over the full-exposure graph, robust to within-user exposure
dependence to second order in the imputation error, and exposure-invariant in
expectation; in simulation, cross-fitting the nuisances on the same log does
not preserve these guarantees. On three benchmarks with unbiased test data, training-free DR
propagation is competitive with tuned trained and closed-form baselines, while
the correction itself leaves top-$K$ accuracy and rankings unchanged; its
contribution is at the level of scores. Future work includes
private nuisance models, certificates that remain informative for large
catalogues, guarantees under adaptive logging, and combining the correction with
learned propagation weights.
